\section{Virtual Memory}
\label{sec:vm}
\sources{Lecture ``Virtual Memory'' (\nolinkurl{05-virtual-memory-lecture.pdf});
Group~2 ``Virtual Memory'' (\nolinkurl{05-virtual-memory-group2.pdf}); H\&P 5th
ed.\ App.~B (\S B.4--B.5) and \S2.4.}

% =============================================================================
\subsection{Why Virtual Memory}
\label{sec:vm-why}

Computers run many processes at once, each wanting a full address space, but
they must share a smaller physical memory. Without virtual memory, programmers
split large programs into \emph{overlays} by hand, programs could overwrite each
other or the OS, free memory became scattered (fragmentation), and whole programs
were loaded up front.\HPA[42]{B}{\S B.4}{40}

\begin{definition}[Virtual memory]
\textbf{Virtual memory} (VM) lets main memory act as a cache for secondary
storage. The processor issues \textbf{virtual} (logical) addresses; hardware (the
MMU) and the OS translate them to \textbf{physical} addresses in main memory or
on disk (\emph{memory mapping}, \emph{address translation}). Each process sees a
large, private address space.
\end{definition}

\paragraph{Benefits.}
\begin{itemize}
  \item \textbf{Larger space}: a program may be larger than physical memory; cold
  pages live on disk (demand paging loads a page on its first access).
  \item \textbf{Isolation and protection}: one page table per process, so the same
  virtual address in two programs maps to different physical frames.
  \item \textbf{Relocation}: a program runs anywhere in physical memory by changing
  the mapping only.
  \item \textbf{Sharing}: several page tables can map the same frame (shared code,
  libraries).
  \item \textbf{Non-contiguous placement}: consecutive virtual pages (e.g.\
  \texttt{A[n]} and \texttt{A[n+1]} across a page boundary) need not be adjacent in
  physical memory.
\end{itemize}
\textbf{Costs}: page tables use memory, translation adds latency (hidden by the
TLB), page faults stall for millions of cycles, and the MMU and TLB cost hardware.

\begin{remark}
\slidefix{} Group~2 lists page-fault stalls as ``disk reads are 1000$\times$
slower''. With a 100~ns memory and a 5--10~ms disk the ratio is about $10^5$, as
their own AMAT example (10~ms vs.\ 100~ns) shows.
\end{remark}

\begin{center}\begin{minipage}{\linewidth}\centering
\begin{tikzpicture}[font=\small,x=1cm,y=0.55cm]
  \node[font=\small\bfseries] at (0.6,4.6) {Virtual};
  \foreach \l/\y/\a in {A/3/0,B/2/4K,C/1/8K,D/0/12K} {
    \draw (0,\y) rectangle ++(1.2,1) node[midway] {\l};
    \node[anchor=east,font=\scriptsize] at (-0.1,\y+0.5) {\a};}
  \node[font=\small\bfseries] at (4.6,4.6) {Physical};
  \foreach \y/\a/\l in {3/0/,2/4K/C,1/8K/,0/12K/,-1/16K/A,-2/20K/,-3/24K/B,-4/28K/} {
    \draw (4,\y) rectangle ++(1.2,1) node[midway] {\l};
    \node[anchor=west,font=\scriptsize] at (5.3,\y+0.5) {\a};}
  \draw[fill=black!10] (7,0) rectangle ++(1.4,1) node[midway] {D};
  \node[font=\small\bfseries] at (7.7,4.6) {Disk};
  \draw[-{Latex}] (1.2,3.5) -- (4,-0.5);
  \draw[-{Latex}] (1.2,2.5) -- (4,-2.5);
  \draw[-{Latex}] (1.2,1.5) -- (4,2.5);
  \draw[-{Latex},dashed] (0.6,0) -- (0.6,-5) -- (7.7,-5) -- (7.7,0);
  \node[font=\scriptsize\itshape,fill=white,inner sep=1pt] at (2.6,-5) {page on disk};
\end{tikzpicture}
\diagramsource{\nolinkurl{05-virtual-memory-lecture.pdf}, handout p.~3 (``Example of virtual memory''); cf.\ H\&P Fig.~B.19, \hppage{B}{41}}
\end{minipage}
\end{center}
\noindent A program of four 4 KB pages: A, B, C are in memory (anywhere), D is on
disk.\HPA{B}{\S B.4, Figure~B.19}{41}


% =============================================================================
\subsection{Paging Vocabulary}
\label{sec:vm-terms}

\begin{definition}[Pages and frames]
A \textbf{(virtual) page} is a fixed-size block of virtual memory; a
\textbf{(physical) frame} is a block of physical memory of the same size (typically
4 KiB; 512 bytes to 16 MB). Any page fits in any frame. The \textbf{page table}
(one per process, indexed by the virtual page number) maps pages to frames; each
\textbf{page table entry} (PTE) holds the frame number and control bits.
\end{definition}

\begin{center}
\small
\begin{tabular}{@{}lP{10.6cm}@{}}
\toprule
Term & Meaning \\
\midrule
VPN & virtual page number: the upper bits of the virtual address \\
PFN / PPN & physical frame (page) number \\
PTE & page table entry: PFN, valid (present), protection (R/W/X, user/supervisor), dirty, reference bits \\
MMU & memory management unit: hardware that translates every access and checks permissions \\
TLB & translation lookaside buffer: a cache of recent translations \\
ASID & address space identifier: tags TLB entries per process, avoiding flushes on a context switch \\
\bottomrule
\end{tabular}
\end{center}

\begin{center}
\small
\begin{tabular}{@{}lll@{}}
\toprule
Cache term & VM term & \\
\midrule
block (line) & page or segment & unit of transfer \\
cache miss & page fault (address fault) & data not in the upper level \\
hit time 1--2 cycles & 40--100 cycles & \\
miss penalty 8--100 cycles & 700{,}000--6{,}000{,}000 cycles & \\
miss rate 0.5--10\% & 0.00001--0.001\% & \\
block 12--128 bytes & page 4096--65{,}536 bytes & \\
size 0.016--1 MB & 4 MB--4 GB & \\
\bottomrule
\end{tabular}
\end{center}
Everything except the miss rate is much worse for VM. Further differences:
cache replacement is done by hardware, page replacement by the OS (so it can
afford to be smart); the processor address size sets the VM size, while the cache
size is independent of it; disk also holds the file system.\HPA[43]{B}{\S B.4, Figure~B.20}{42}

\begin{definition}[Fragmentation]
\textbf{Internal}: fixed-size allocation wastes the unused end of the last page
(on average half a page per region; about 1.5 pages per process with text, heap
and stack). \textbf{External}: variable-size allocation leaves free memory split
into holes too small to use, although their total is large enough. Paging has
only internal fragmentation; segmentation has external
fragmentation.\HPA{B}{\S B.4}{47}
\end{definition}

% =============================================================================
\subsection{Address Translation and Page Tables}
\label{sec:vm-translation}

\begin{center}\begin{minipage}{\linewidth}\centering
\begin{tikzpicture}[font=\small]
  \node[draw,minimum width=4.4cm,minimum height=0.6cm] (vpn) {Virtual page number};
  \node[draw,minimum width=3cm,minimum height=0.6cm,right=0pt of vpn] (off) {Page offset};
  \node[draw,minimum width=2.6cm,minimum height=1.6cm,below=0.9cm of vpn,align=center] (pt) {Page table\\(V, prot., D, R,\\PFN)};
  \node[draw,minimum width=4.4cm,minimum height=0.6cm,below=0.9cm of pt] (pfn) {Physical page number};
  \node[draw,minimum width=3cm,minimum height=0.6cm,right=0pt of pfn] (off2) {Page offset};
  \node[left=0.6cm of pt,align=center,font=\scriptsize] (ptbr) {page-table\\base register};
  \draw[-{Latex}] (vpn) -- node[right,font=\scriptsize]{index} (pt);
  \draw[-{Latex}] (ptbr) -- (pt);
  \draw[-{Latex}] (pt) -- (pfn);
  \draw[-{Latex}] (off.south) -- (off2.north);
  \node[right=2pt of off,font=\scriptsize] {$\log_2$ page size bits};
\end{tikzpicture}
\diagramsource{\nolinkurl{05-virtual-memory-lecture.pdf}, handout p.~2 (``VM: Address Translation''); Group~2 slide 13 (``Page Mapping Translation Flow''); cf.\ H\&P Fig.~B.23, \hppage{B}{45}}
\end{minipage}
\end{center}

\begin{prop}[Translation]
With page size $P$:
\[
  \text{VPN}=\left\lfloor\frac{\text{VA}}{P}\right\rfloor,\qquad
  \text{offset}=\text{VA}\bmod P,\qquad
  \text{PA}=\text{PFN}\times P+\text{offset}.
\]
The offset is not translated; only the page number changes. For paging the offset
is \emph{concatenated} to the frame number; for segmentation it is
\emph{added} to the segment base.\HPA{B}{\S B.4, Figure~B.23}{44}
\end{prop}

\begin{eg}[Hex and decimal translation (Group 2)]
32-bit VA, 4 KiB pages (12-bit offset $=$ the last 3 hex digits). VA
\texttt{0x00403A7C}: VPN \texttt{0x00403}, offset \texttt{0xA7C}; the PTE gives PFN
\texttt{0x1F2}, so PA $=$ \texttt{0x1F2A7C}.

Page size 1024: VA 3500 has VPN $\lfloor3500/1024\rfloor=3$ and offset
$3500-3072=428$; page 3 maps to frame 7, so PA $=7\times1024+428=7596$.
\end{eg}

\begin{prop}[Page table size]
A flat table has one PTE per virtual page:
size $=2^{\text{VA bits}-\log_2P}\times$ PTE size, per process. For a 32-bit VA,
4 KB pages and 4-byte PTEs: $2^{20}\times4=4$ MB per process (400 MB for 100
processes), mostly mapping nothing.\HPA{B}{\S B.4}{44}
\end{prop}
\begin{itemize}
  \item \textbf{Multilevel page tables}: split the VPN; allocate second-level
  tables only for used regions. 32-bit: $10+10+12$ (each level fits in a 4 KB page
  of 1024 four-byte PTEs); x86-64: $9+9+9+9+12$ (48-bit VA, 256 TiB); a TLB miss
  costs one memory access per level.
  \item \textbf{Inverted page table}: one entry per physical frame, found by hashing
  the VPN; its size follows physical memory. A 512 MB memory with 4 KB pages needs
  only 1 MB (8 bytes per entry, including the virtual address) instead of 4 MB per
  process.\HPA[45]{B}{\S B.4}{44}
  \item On x86 the register CR3 holds the root of the active page table and is
  switched on every context switch.
\end{itemize}

% =============================================================================
\subsection{Page Faults}
\label{sec:vm-faults}

\begin{definition}[Page fault]
A \textbf{page fault} is the exception raised when the PTE of an accessed page is
not valid (the page is not in memory). It is the VM analogue of a cache miss, but
handled in \emph{software}: the processor switches to another task during the
disk access.\HPA{B}{\S B.1}{3}
\end{definition}
\begin{enumerate}
  \item The MMU finds a missing or invalid mapping and traps to the OS.
  \item The OS checks the access is legal; an illegal one is rejected (``segmentation
  fault'').
  \item It takes a free frame or chooses a victim; a \textbf{dirty} victim is written
  back first.
  \item It reads the page from disk (or zero-fills a new page).
  \item It updates the PTE and restarts the faulting instruction.
\end{enumerate}
Pinned pages (the fault handler, DMA buffers) never page out. \textbf{Thrashing}:
the working sets of the running processes exceed memory, so the system spends
more time paging than executing; add memory or run fewer processes.

\begin{prop}[Memory access time with page faults]
With hit rate $h$, memory time $T_{mem}$ and disk (fault) time $T_{disk}$, the
lecturer's slide writes $\AMAT=T_{mem}+(1-h)\,T_{disk}$; Group~2 writes
$h\,T_{mem}+(1-h)\,T_{disk}$. They differ by $(1-h)\,T_{mem}$, negligible because
$1-h$ is tiny.
\end{prop}
\begin{eg}[How bad is 99.9\%?]
$T_{mem}=100$ ns, $T_{disk}=10$ ms $=10^7$ ns, $h=99.9\%$:
$0.999\times100+0.001\times10^7\approx10{,}100$ ns, about $101\times$ slower than
memory. With the lecturer's numbers ($T_{disk}=25$ ms) it is $25{,}100$ ns. The
hit rate must be almost perfect.
\end{eg}

% =============================================================================
\subsection{The Four Questions for Virtual Memory}
\label{sec:vm-questions}

\begin{itemize}
  \item \textbf{Q1 Placement}: anywhere in main memory (\textbf{fully associative}):
  the huge miss penalty makes a lower miss rate worth a complex placement.
  \item \textbf{Q2 Identification}: the page table indexed by the VPN (or an inverted
  table), sped up by the TLB.
  \item \textbf{Q3 Replacement}: approximately \textbf{LRU}, by the OS. Hardware
  provides a \emph{use} (reference) bit, set on access; the OS clears the bits
  periodically to see which pages were used recently.
  \item \textbf{Q4 Writes}: always \textbf{write-back} with a dirty bit; nobody
  writes through to disk on every store.
\end{itemize}
\HPA[46]{B}{\S B.4}{44}

\subsubsection{Page replacement algorithms}
\label{sec:vm-replacement}
\begin{center}
\small
\begin{tabular}{@{}lP{5.4cm}P{6.2cm}@{}}
\toprule
Policy & Evicts & Notes \\
\midrule
OPT (optimal) & the page next used furthest in the future & needs the future: a benchmark only \\
FIFO & the page loaded earliest & simple; may evict a busy page; Belady's anomaly \\
LRU & the page unused for the longest time & uses temporal locality; exact tracking is costly \\
LFU & the least frequently used page & keeps once-popular stale pages \\
Clock & first page with reference bit 0 (clearing 1s on the way) & cheap LRU approximation \\
\bottomrule
\end{tabular}
\end{center}

\begin{eg}[FIFO vs.\ LRU (Group 2)]
References \texttt{7 0 1 2 0 3 0 4}, 3 frames. FIFO: 7, 0, 1 fault; 2 evicts 7; 0
hits; 3 evicts 0 (though just used); 0 evicts 1; 4 evicts 2: \textbf{7 faults}. LRU:
2 evicts 7; 0 hits; 3 evicts 1; 0 hits; 4 evicts 2: \textbf{6 faults}.
\end{eg}

\begin{eg}[A string where FIFO beats LRU (lecturer's review)]
3 frames, \texttt{1 2 1 1 3 2 1 4 3 1 1 2 4 1 5 6 2 1}; F $=$ fault (victim).
\begin{center}
\scriptsize
\setlength{\tabcolsep}{2.2pt}
\begin{tabular}{@{}l*{18}{c}@{}}
\toprule
Ref & 1 & 2 & 1 & 1 & 3 & 2 & 1 & 4 & 3 & 1 & 1 & 2 & 4 & 1 & 5 & 6 & 2 & 1 \\
\midrule
FIFO & F & F & & & F & & & F(1) & & F(2) & & F(3) & & & F(4) & F(1) & & F(2) \\
LRU & F & F & & & F & & & F(3) & F(2) & & & F(4) & F(3) & & F(2) & F(4) & F(1) & F(5) \\
OPT & F & F & & & F & & & F(2) & & & & F(3) & & & F(4) & F(5) & & \\
\bottomrule
\end{tabular}
\end{center}
FIFO 9, LRU 11, OPT 7 faults. At reference 4 LRU evicts 3 (used longest ago) and
needs it next; FIFO evicts 1 (loaded first). LRU is usually better, but not on every
string.
\end{eg}

\begin{remark}
\textbf{Belady's anomaly}: with FIFO, more frames can give \emph{more} faults:
\texttt{1 2 3 4 1 2 5 1 2 3 4 5} gives 9 faults with 3 frames and 10 with 4. LRU
and OPT are \emph{stack algorithms} (the pages kept with $n$ frames are always
kept with $n+1$) and never show it.
\end{remark}

% =============================================================================
\subsection{The TLB}
\label{sec:tlb}

\begin{definition}[Translation lookaside buffer]
Page tables live in memory, so every access would need two memory accesses
(PTE, then data). The \textbf{TLB} (translation buffer) caches recent
translations: the tag holds part of the VPN, the data holds the PFN, protection,
valid, use and dirty bits. Locality makes a small TLB effective.\HPA{B}{\S B.4}{46}
\end{definition}
\begin{itemize}
  \item \textbf{TLB hit}: PA $=$ PFN $\Vert$ offset at once, no page-table access.
  \item \textbf{TLB miss}: walk the page table (hardware or OS); if the page is
  present, load the PTE into the TLB; if not, \textbf{page fault}. A TLB miss is
  not a page fault.
  \item \textbf{Opteron data TLB}: 40 entries, fully associative; the VA is sent to
  all tags while the access type is checked against the protection bits; a 40:1
  multiplexer selects the 28-bit PFN, joined to the 12-bit offset (40-bit PA).
  \item On a context switch: flush the TLB, or tag entries with an ASID/PCID.
  \item \textbf{TLB reach} $=$ entries $\times$ page size; large pages map more
  memory with the same entries.\HPA{B}{\S B.4, Figure~B.24}{47}
\end{itemize}

\begin{prop}[Effective access time with a TLB]
TLB time $t$, memory time $m$, TLB hit rate $h$, one-level page table:
\[
  \text{EAT}=h\,(t+m)+(1-h)\,(t+2m).
\]
A $k$-level table costs $t+(k+1)m$ on a miss.
\end{prop}
\begin{eg}[EAT (Group 2)]
$t=20$, $m=100$ ns, $h=98\%$: $0.98\times120+0.02\times220=122$ ns, a 22\%
overhead over 100 ns. With $t=10$, $m=80$, $h=95\%$:
$0.95\times90+0.05\times170=94$ ns.
\end{eg}

\begin{center}
\small
\begin{tabular}{@{}lllp{6.2cm}@{}}
\toprule
TLB & Page table & Cache & Possible? \\
\midrule
hit & -- & hit or miss & yes (the page table is not accessed) \\
miss & hit & hit or miss & yes \\
miss & miss (fault) & miss & yes \\
hit & miss & any & no: a page in the TLB is in memory \\
miss & miss & hit & no: data cannot be cached if its page is not in memory \\
\bottomrule
\end{tabular}
\end{center}

\paragraph{Selecting a page size.} Larger pages: smaller page tables, larger
VIPT caches, more efficient disk transfers, more memory mapped per TLB entry (fewer
TLB misses), hence multiple page sizes in modern processors. Smaller pages: less
internal fragmentation and faster process start-up.\HPA[47]{B}{\S B.4}{46}

\subsubsection{From virtual address to cache}
\label{sec:vm-cache}
A \textbf{virtually indexed, physically tagged} (VIPT) L1 cache uses page-offset
bits as its index, so the cache is read while the TLB translates; the physical tag
is then compared. This needs index $+$ block offset $\le$ page offset, i.e.\
$C/a\le$ page size; to grow the L1, raise the associativity.\HPA{B}{\S B.4, Figure~B.25}{48}
\begin{eg}[A whole hierarchy]
64-bit VA, 41-bit PA, 8 KB pages, direct-mapped 256-entry TLB, direct-mapped 8 KB
L1 and 4 MB L2, 64-byte blocks.
\begin{center}
\small
\begin{tabular}{@{}lr@{}}
\toprule
VA $=$ VPN $+$ page offset & $51+13$ \\
VPN $=$ TLB tag $+$ TLB index & $43+8$ \\
L1: index $+$ block offset $=$ page offset & $7+6=13$ \\
L1 tag $=$ physical page number & $41-13=28$ \\
L2: tag $+$ index $+$ block offset & $19+16+6=41$ \\
\bottomrule
\end{tabular}
\end{center}
\end{eg}

% =============================================================================
\subsection{Segmentation}
\label{sec:segmentation}

\begin{definition}[Segmentation]
Memory is divided into \textbf{variable-size segments} that follow the program's
logical structure (code, data, heap, stack) and may be visible to the programmer.
An address is (segment number, offset); the segment table holds each segment's
\textbf{base}, \textbf{limit} and protection:
\[
  \text{if offset}<\text{limit}:\ \text{PA}=\text{base}+\text{offset};\qquad
  \text{else trap.}
\]
\HPA[44]{B}{\S B.4, Figures~B.21--B.22}{43}
\end{definition}

\begin{eg}[Segment translation (Group 2)]
Segment 2 (stack): base 4300, limit 400. (2, 53): $53<400$, PA $=4353$. (2, 500):
$500\ge400$, trap, no address formed. Segment 1: base 6300, limit 400: (1, 399)
$\to6699$; (1, 400) traps: valid offsets are $0,\dots,\text{limit}-1$.
\end{eg}

\begin{center}
\small
\begin{tabular}{@{}lP{5.8cm}P{5.8cm}@{}}
\toprule
& Paging & Segmentation \\
\midrule
Unit & fixed-size page & variable-size segment \\
Words per address & one & two (segment, offset) \\
Visible to programmer & no & may be \\
Replacing a block & trivial (all the same size) & hard (find a contiguous free region) \\
Fragmentation & internal & external (holes; needs compaction; first/best/worst fit) \\
Disk traffic & efficient (tune the page size) & small segments may move few bytes \\
Protection, sharing & fine-grained, per page & natural, per logical region \\
\bottomrule
\end{tabular}
\end{center}
\begin{itemize}
  \item \textbf{Paged segments} (segmented paging): each segment is a whole number
  of pages with its own page table. Replacement is easy, yet a segment is a logical
  unit. Up to three memory reads uncached (segment table, page table, data);
  Multics, IBM System/38.
  \item Because of the replacement problem few computers use pure segmentation.
  \item \textbf{x86}: real mode PA $=$ segment $\times16+$ offset (20 bits, 1 MiB,
  no protection; \texttt{06EF:9234} $\to$ \texttt{0x10124}); protected mode
  (80286) uses selectors into descriptor tables with base, limit and privilege
  rings 0--3; from the 80386, paging under segmentation, and today's OSes use a
  flat model.
\end{itemize}

% =============================================================================
\subsection{Protection}
\label{sec:protection}

\begin{itemize}
  \item \textbf{Isolation}: each process has its own address space and page table;
  a process cannot reach memory it has no mapping for.
  \item \textbf{Permissions}: R/W/X per page (code RX, data RW, constants R) or per
  segment; writing read-only memory or executing data traps.
  \item \textbf{Privilege}: user vs.\ kernel (supervisor) mode; U/S bit per page.
\end{itemize}
At a minimum the architecture must (1) provide at least \textbf{two modes},
user and supervisor (kernel); (2) provide processor state a user process can use
but not write (the mode bit, the exception enable bit, memory protection
information); (3) provide ways to switch modes: a \emph{system call} to supervisor
mode and a return that restores user mode; (4) limit memory accesses so a process
need not be swapped to disk on a context switch.\HP{\S2.4}{106} Processes are
protected by having their own page tables, which user programs cannot modify;
protection can be extended with rings (80x86) or capabilities. The check happens
during translation, in the TLB, on every access.\HPA{B}{\S B.5}{50}

% =============================================================================
\subsection{Worked Problems}
\label{sec:vm-problems}

\begin{eg}[Address split (Group 2)]
16-bit addresses, 256-byte pages: offset $\log_2256=8$ bits, VPN $16-8=8$ bits, so
$2^8=256$ virtual pages. 2 KiB pages, VA 5000, page 2 $\to$ frame 9: VPN
$\lfloor5000/2048\rfloor=2$, offset $5000-4096=904$, PA $=9\times2048+904=19{,}336$.
\end{eg}

% =============================================================================
\subsection{Textbook Exercises}
\label{sec:vm-tbx}

\begin{tbexercise}[B]{B.12}{64}{\S B.4 (TLB, page tables), \hppages{B}{44}{47}}
A program runs on a computer with a four-entry fully associative TLB:
\begin{center}
\small
\begin{tabular}{@{}ccc@{\qquad}ccc@{\qquad}ccc@{}}
\toprule
VP\# & PP\# & Valid & VP\# & PP\# & Valid & VP\# & PP\# & Valid \\
\midrule
5 & 30 & 1 & 7 & 1 & 0 & 10 & 10 & 1 \\
15 & 25 & 1 & & & & & & \\
\bottomrule
\end{tabular}
\end{center}
Page table (virtual page: physical page, present): 0: 3 Y; 1: 7 N; 2: 6 N; 3: 5 Y;
4: 14 Y; 5: 30 Y; 6: 26 Y; 7: 11 Y; 8: 13 N; 9: 18 N; 10: 10 Y; 11: 56 Y; 12: 110 Y;
13: 33 Y; 14: 12 N; 15: 25 Y. For the virtual page trace 1, 5, 9, 14, 10, 6, 15,
12, 7, 2, say for each access whether the TLB hits or misses and, if the page
table is accessed, whether it is a page hit or a page fault.
\end{tbexercise}

\begin{tbexercise}[B]{B.13}{66}{\S B.4, \hppages{B}{46}{47}}
Some systems handle TLB misses in software (as an exception), others in
hardware. (a) What are the trade-offs? (b) Is software handling always slower?
(c) Are there page-table structures easy in software but hard in hardware, or the
reverse? (d) Why are TLB miss rates of floating-point programs generally higher
than those of integer programs?
\end{tbexercise}

\begin{tbexercise}[B]{B.15}{66}{\S B.5, \hppages{B}{49}{56}}
A scheme like HP/PA's stores a protection ID (key) in each PTE; on every access the
CPU compares it with four protection-ID registers (writable only in supervisor
mode) and checks the access rights. (a) How could a process hold more than four
valid protection IDs at a time (say ten)? Propose a mechanism, perhaps with
software help. (The book has further parts on the same scheme.)
\end{tbexercise}
